\documentclass[11pt,letterpaper]{article}

% Essential Packages
\usepackage[margin=1in]{geometry}
\usepackage{amsmath,amssymb}
\usepackage{graphicx}
\usepackage{circuitikz}
\usepackage{tikz}
\usetikzlibrary{shapes,arrows,positioning,calc}
\usepackage{booktabs}
\usepackage{multirow}
\usepackage{xcolor}
\usepackage{fancyhdr}
\usepackage{tcolorbox}
\usepackage{enumitem}
\usepackage{float}
\usepackage{listings}
\usepackage{hyperref}

% Colors
\definecolor{nlcblue}{RGB}{0,51,102}
\definecolor{alertred}{RGB}{204,0,0}
\definecolor{safgreen}{RGB}{0,128,0}

% Header/Footer
\pagestyle{fancy}
\fancyhead[L]{\textbf{NLC STEM Club UAV Project}}
\fancyhead[R]{\textbf{Project Overview}}
\fancyfoot[C]{\thepage}

% Title
\title{\Huge\textbf{Project Overview}\\
\Large 3D-Printed Solar Floatplane UAV\\
\large NLC STEM Club - Electrical Systems Team}
\author{February 2026}
\date{}

\begin{document}

\maketitle
\thispagestyle{empty}
\newpage
NLC STEM Club: Solar-Powered Environmental Monitoring UAV\par

This project is a high-impact, interdisciplinary effort to design, build, and fly a solar-powered drone for environmental research. We will work as a single team, divided into specialized groups, to achieve a common goal. This project will serve as a powerful portfolio piece for every member, demonstrating your ability to execute a complex engineering mission from start to finish.\par

The Project's Main Goal and Scope\par

Our mission is to create a fully-functional, hybrid solar-powered drone that can collect critical environmental data. This project is a tangible example of a modern engineering challenge, combining mechanical design, electrical systems, and software.\par

Primary Objective: Build an operational UAV that uses solar energy to extend its flight time beyond a standard battery.\par

Data Collection: The drone will collect real-time data on air quality, temperature, and humidity.\par

Final Product: The project will culminate in a final report that documents the entire process and presents our findings.\par

Predicted Drone Parameters\par

Based on our budget and a realistic timeline, we have defined the following parameters for our drone. These are not final and may be adjusted as we progress.\par

Weight: Approximately 4-6 pounds, including all electronics and the payload. We must keep the weight low to maximize flight time.\par

Size: The drone will be a quadcopter with a frame size of around 250mm (measured diagonally from motor to motor). This is a standard and stable size.\par

Flight Time: We predict an extended flight time of 15-25 minutes thanks to the solar-charging system. This is a significant improvement for a drone of this size.\par

\subsection{Sensors: The drone will be equipped with a suite of environmental sensors to measure}

Air Quality (VOC, CO₂)\par

Temperature and Humidity\par

The Teams and Their Roles\par

This project is divided into four distinct teams, each led by an officer. Every team will work in parallel to ensure we meet our goals.\par

Mechanical \& Structural Team: The Builders. They will design and 3D print the drone's body, ensuring it is both strong and lightweight enough to fly. They are the ones who turn the design into a physical reality.\par

Electronics \& Power Team: The Integrators. They will handle all the wiring and electrical systems, including the critical solar power system. They are responsible for bringing the drone to life.\par

Software \& Data Team: The Brains. They will write the code that makes the drone fly and collect data. They will also analyze all the collected data and turn it into a final report.\par

Documentation \& Reporting Team: The Storytellers. They are responsible for recording every step of the project, from design to flight. They will write the final research paper and create the video that showcases our work.\par

Project Timeline\par

This project will follow a 10-week timeline. All teams will work on their respective tasks simultaneously to ensure we can meet our final deadline.\par

1: Design \& Planning (1-3)\par

All teams finalize their designs and identify necessary components.\par

2: Fabrication \& Assembly (4-6)\par

The physical drone is built, and all electronic components are wired and tested.\par

3: Integration \& Testing (7-9)\par

All systems are integrated (hardware and software) and tested together.\par

4: Final Report \& Flight (10)\par

We fly the drone, analyze the final data, and complete the project's documentation.\par

What We Need From You\par

This project requires everyone's effort. We want members who are willing to learn, communicate, and contribute to their team's goals. Whether you're a seasoned builder or a complete beginner, there's a place for you on a team. We can teach you everything you need to know, but you must be willing to show up and contribute.\par

Your contribution is what will make this project a success and a memorable experience.\par

\end{document}
